% CP-VI-0045
% target_chapter: VI/38 — Projective Morphisms and Closed Embeddings
% source_unit_id: IMP-DL-0A9B1F0A75-P002
% source: Downloads/theory-of-algebraic-geometry-9.tex L822-L1065
% solution_provenance: SOURCE_EXPLICIT_SOLUTION

\subsection*{Problem statement}

Let
\[
\varphi:\mathbf{P}^n\longrightarrow \mathbf{P}^m
\]
be a morphism. Show that either
\[
\varphi(\mathbf{P}^n)
\]
is a point, or else
\[
m\geq n
\qquad\text{and}\qquad
\dim \varphi(\mathbf{P}^n)=n.
\]

The suggested strategy is first to prove that if \(m<n\), then every morphism
\[
\mathbf{P}^n\longrightarrow \mathbf{P}^m
\]
is constant, and then to finish by using projections from linear subspaces.

\subsection{Morphisms to lower-dimensional projective spaces}

\begin{lemma}
	Let
	\[
	\psi:\mathbf{P}^n\longrightarrow \mathbf{P}^r
	\]
	be a morphism with
	\[
	r<n.
	\]
	Then \(\psi\) is constant.
\end{lemma}

\begin{proof}
	Since
	\[
	\operatorname{Pic}(\mathbf{P}^n)\cong \mathbb{Z},
	\]
	there exists an integer \(d\geq 0\) such that
	\[
	\psi^*\mathcal{O}_{\mathbf{P}^r}(1)\cong \mathcal{O}_{\mathbf{P}^n}(d).
	\]
	
	We distinguish two cases.
	
	\subsubsection*{Case 1: \(d=0\)}
	
	If
	\[
	\psi^*\mathcal{O}_{\mathbf{P}^r}(1)\cong \mathcal{O}_{\mathbf{P}^n},
	\]
	then the pullbacks of the homogeneous coordinate sections on \(\mathbf{P}^r\) are global sections of the trivial line bundle on \(\mathbf{P}^n\).
	
	But
	\[
	\Gamma(\mathbf{P}^n,\mathcal{O}_{\mathbf{P}^n})=k.
	\]
	
	Thus these pullback sections are constants. Hence all points of \(\mathbf{P}^n\) are sent to the same point of \(\mathbf{P}^r\), and \(\psi\) is constant.
	
	\subsubsection*{Case 2: \(d>0\)}
	
	If \(d>0\), then \(\psi\) is represented by \(r+1\) homogeneous polynomials of the same degree \(d\):
	\[
	F_0,\dots,F_r\in k[X_0,\dots,X_n],
	\]
	with no common zero in \(\mathbf{P}^n\).
	
	Indeed,
	\[
	\psi([X_0:\cdots:X_n])
	=
	[F_0(X):\cdots:F_r(X)].
	\]
	
	Since \(r<n\), we have
	\[
	r+1\leq n.
	\]
	
	Consider the homogeneous ideal
	\[
	I=(F_0,\dots,F_r)\subseteq k[X_0,\dots,X_n].
	\]
	
	By Krull's height theorem,
	\[
	\operatorname{ht}(I)\leq r+1\leq n.
	\]
	
	On the other hand, if the homogeneous polynomials \(F_0,\dots,F_r\) had no common zero in projective space, then the radical of \(I\) would contain the irrelevant ideal
	\[
	(X_0,\dots,X_n),
	\]
	whose height is
	\[
	n+1.
	\]
	
	That would force
	\[
	\operatorname{ht}(I)=n+1,
	\]
	contradicting
	\[
	\operatorname{ht}(I)\leq n.
	\]
	
	Therefore the polynomials \(F_0,\dots,F_r\) must have a common zero in \(\mathbf{P}^n\). But then they do not define a morphism on all of \(\mathbf{P}^n\), a contradiction.
	
	Hence the case \(d>0\) is impossible. Therefore \(d=0\), and \(\psi\) is constant.
\end{proof}

Thus we have proved:
\[
\boxed{
	r<n
	\quad\Longrightarrow\quad
	\text{every morphism }\mathbf{P}^n\to \mathbf{P}^r\text{ is constant.}
}
\]

\subsection{Dimension of the image of a nonconstant morphism}

\textbf{Solution.}

Let
\[
Y=\varphi(\mathbf{P}^n)\subseteq \mathbf{P}^m.
\]

Since \(\mathbf{P}^n\) is irreducible and \(\varphi\) is continuous, \(Y\) is irreducible. Put
\[
r=\dim Y.
\]

Because \(Y\) is the image of an \(n\)-dimensional variety,
\[
r\leq n.
\]

We shall show that if \(Y\) is not a point, then in fact
\[
r=n.
\]

Suppose, for contradiction, that
\[
0<r<n.
\]

After extending the base field if necessary, we may choose a sufficiently general linear subspace
\[
\Lambda\subseteq \mathbf{P}^m
\]
of dimension
\[
m-r-1
\]
such that
\[
\Lambda\cap Y=\varnothing.
\]

Projection away from \(\Lambda\) defines a morphism
\[
\pi_{\Lambda}:Y\longrightarrow \mathbf{P}^r.
\]

For a sufficiently general choice of \(\Lambda\), this linear projection is finite onto its image. Since
\[
\dim Y=r>0,
\]
the morphism \(\pi_{\Lambda}\) is nonconstant.

Now consider the composite morphism
\[
\mathbf{P}^n
\xrightarrow{\;\varphi\;}
Y
\xrightarrow{\;\pi_{\Lambda}\;}
\mathbf{P}^r.
\]

This is a nonconstant morphism
\[
\mathbf{P}^n\longrightarrow \mathbf{P}^r.
\]

But
\[
r<n,
\]
so by the lemma every such morphism must be constant. This is a contradiction.

Therefore no intermediate dimension
\[
0<r<n
\]
is possible. Hence
\[
r=0
\qquad\text{or}\qquad
r=n.
\]

If
\[
r=0,
\]
then the irreducible projective variety \(Y\) is a single point. Thus
\[
\varphi(\mathbf{P}^n)
\]
is a point.

If
\[
r=n,
\]
then \(Y\subseteq \mathbf{P}^m\) has dimension \(n\). Since every closed subset of \(\mathbf{P}^m\) has dimension at most \(m\), we must have
\[
n\leq m.
\]

Therefore
\[
\boxed{
	\text{either }\varphi(\mathbf{P}^n)\text{ is a point, or else }
	m\geq n
	\text{ and }
	\dim\varphi(\mathbf{P}^n)=n.
}
\]

\begin{remark}
	The result says that projective space cannot be mapped nontrivially onto a strictly lower-dimensional projective variety by a morphism. A nonconstant morphism from \(\mathbf{P}^n\) must preserve the full dimension of \(\mathbf{P}^n\).
\end{remark}

% ============================================================
